In this appendix we show the computation of the parameters of the hypergeometric functions and their relation with the rotation parameters.


\subsection{Consistency Conditions of the Monodromy Matrices}

In the main text we set
\begin{equation}
  D~
  \rM_{\vb{\infty}}~
  D^{-1}
  =
  e^{-2\pi i \delta_{\vb{\infty}}}\,
  \cL(\vb{n}_{\vb{\infty}}),
\end{equation}
where $\cL(\vb{n}_{\vb{\infty}}) \in \SU{2}$.
The previous equation implies
\begin{equation}
  \left( D\, \rM_{\vb{\infty}}\, D^{-1} \right)^\dagger
  =
  \left( D\, \rM_{\vb{\infty}}\, D^{-1} \right)^{-1},
\end{equation}
which can be rewritten as
\begin{equation}
  \widetilde{\rM}_{\vb{\infty}}^{-1}~
  \cC^{\dagger}\, D^{\dagger}\, D\, \cC
  =
  \cC^{\dagger}\, D^{\dagger}\, D\, \cC~
  \widetilde{\rM}_{\vb{\infty}}^{-1}.
\end{equation}
As $\widetilde{\rM}_{\vb{\infty}}$ is a generic diagonal matrix, the previous equation implies that the off-diagonal elements of $\cC^{\dagger}\, D^{\dagger}\, D\, \cC$ must vanish.
We therefore have
\begin{equation}
  \begin{split}
    \abs{K}^{-2}
    & =
    -\frac{\cC_{21}\, \cC^*_{22}}{\cC_{11}\, \cC^*_{12}}
    \\
    & =
    -\frac{1}{\pi^4}\,
    \abs{\gfun{a} \gfun{b} \gfun{c-a} \gfun{c-b}}^2 \times
    \\
    & \times
    \sin(\pi a)\, \sin^*(\pi (c-a))\, (\sin(\pi b)\, \sin^*(\pi (c-b)))^*.
  \end{split}
\end{equation}
When $a,\, b,\, c \in \R$ this ultimately means that
\begin{equation}
  \sin(\pi a)\, \sin(\pi (c-a))\, \sin(\pi b)\, \sin(\pi (c-b)) < 0.
  \label{eq:constraint_from_K^2}
\end{equation}
Since the previous equation is invariant under integer shift of any of its parameters, we can consider just the fractional parts $0 \le \{a\},\, \{b\},\, \{c\} < 1$.
In order to have \U{2} monodromies finally requires
\begin{equation}
  0 \le \{b\} < \{c\} < \{a\} < 1
  \qq{or}
  0 \le \{a\} < \{c\} < \{b\} <1.
  \label{eq:K_consistency_condition}
\end{equation}

Should we request \U{1,1} monodromies as in moving rotated branes then we get:
\begin{equation}
  \abs{K}^{-2}
  =
  \frac{\cC_{21}\, \cC^*_{22}}{\cC_{11}\, \cC^*_{12}}.
\end{equation}
This would then imply
\begin{equation}
  0 \le \{c\} < \{a\},\, \{b\} < 1
  \qq{or}
  0 \le \{a\},\, \{b\} < \{c\} < 1.
\end{equation}


\subsection{Fixing the Parameters}

We can finally show in details the computation of the parameters of the basis of hypergeometric functions used in the main text.
The relation between these and the \SU{2} matrices can be computed requiring that the monodromies induced by the choice of the parameters equal the monodromies produced by the rotations of the D-branes.

The monodromy in $\omega_{\bt-1} = 0$ is simpler to compute given that we choose $\cL(\vb{n}_{\vb{0}})$ and $\cR(\widetilde{\vb{m}}_{\vb{0}})$ to be diagonal.
We impose:
\begin{eqnarray}
  \mqty( \dmat{1, e^{-2\pi i c^{(L)}}} )
  & = &
  e^{-2\pi i \delta_{\vb{0}}^{(L)}}\,
  \mqty( \dmat{e^{2\pi i n_{\vb{0}}}, e^{-2\pi i n_{\vb{0}}}} ),
  \\
  \mqty( \dmat{1, e^{-2\pi i c^{(R)}}} )
  & = &
  e^{-2\pi i \delta_{\vb{0}}^{(R)}}\,
  \mqty( \dmat{e^{-2\pi i m_{\vb{0}}}, e^{2\pi i m_{\vb{0}}}} ),
\end{eqnarray}
where $n^3_{\vb{0}} = \norm{\vb{n}_{\vb{0}}} = n_{\vb{0}}$ and $m^3_{\vb{0}} = \norm{\vb{m}_{\vb{0}}} = m_{\vb{0}}$ with $0 \le  n_{\vb{0}},\, m_{\vb{0}} < 1$ due to the conventions \eqref{eq:maximal_torus_left} and \eqref{eq:maximal_torus_right}.
We thus have:
\begin{equation}
  \begin{split}
    \delta_{\vb{0}}^{(L)}
    & =
    n_{\vb{0}} + k_{\delta^{(L)}_{\vb{0}}},
    \qquad
    k_{\delta^{(L)}_{\vb{0}}} \in \Z,
    \\
    c^{(L)}
    & =
    2 n_{\vb{0}} + k_c,
    \qquad
    k_c \in \Z.
  \end{split}
  \label{eq:cL}
\end{equation}
Since the determinant of the right hand side is $e^{-4 \pi i \delta_{\vb{0}}^{(L)}}$, the range of definition of $\delta_{\vb{0}}^{(L)}$ is $\alpha \le \delta_{\vb{0}}^{(L)} \le  \alpha + \frac{1}{2}$.
Given that $0 \le n_{\vb{0}} < \frac{1}{2}$ we simply take $\alpha = 0$ and set $\delta_{\vb{0}}^{(L)} = n_{\vb{0}}$.
Analogous results hold in the right sector.
Furthermore from the third equation in \eqref{eq:parameters_equality_zero} and from the first equation in \eqref{eq:cL} we can restrict:
\begin{equation}
  n_{\vb{0}} + m_{\vb{0}} - A \in \Z.
\end{equation}

We then need to find $3$ equations to determine $a^{(L)}$, $b^{(L)}$ and $\delta^{(L)}_{\vb{\infty}}$.
After that we then fix the remaining factors in $B$ and $\abs{K^{(L)}}$.
The equations follow from~\eqref{eq:parameters_equality_infty}.
The first two equations for $a^{(L)}$, $b^{(L)}$ and $\delta^{(L)}_{\vb{\infty}}$ follow by considering the trace of~\eqref{eq:parameters_equality_infty}:
\begin{equation}
  e^{\pi i ( a^{(L)} + b^{(L)} )} \cos(\pi( a^{(L)} - b^{(L)} ) )
  =
  e^{-2\pi i \delta^{(L)}_{\infty}} \cos(2\pi n_{\vb{\infty}}),
\end{equation}
which is satisfied by:
\begin{equation}
  \begin{split}
    \delta^{(L)}_{\vb{\infty}}
    & =
    -
    \frac{1}{2}(a^{(L)} + b^{(L)})
    +
    \frac{1}{2} k_{\delta^{(L)}_{\vb{\infty}}},
    \qquad
    k_{\delta_{\vb{\infty}}} \in \Z,
    \\
    a^{(L)} - b^{(L)}
    & =
    2\, (-1)^{p^{(L)}}\, n_{\vb{\infty}}
    +
    (-1)^{q^{(L)}}\, k_{\delta^{(L)}_{\vb{\infty}}}
    +
    2\, k'_{a b},
    \qquad
    k'_{ab} \in \Z,
  \end{split}
\end{equation}
where $p^{(L)},\, q^{(L)} \in \left\lbrace 0, 1 \right\rbrace$.
Notice that changing the value of $p^{(L)}$ corresponds to swapping $a$ and $b$: since the hypergeometric function is symmetric in those parameters we can fix $p^{(L)}=0$.
Redefining $k'$ we can always set $q^{(L)}=0$.
We therefore have:
\begin{equation}
  a^{(L)} - b^{(L)}
  =
  2\, n_{\vb{\infty}}
  +
  k_{\delta^{(L)}_{\vb{\infty}}}
  +
  2 k_{ab},
  \qquad
  k_{a b}\in \Z.
  \label{eq:aL-bL}
\end{equation}
The allowed values for $k_{\delta^{(L)}_{\vb{\infty}}}$ follow a construction similar to the monodromy around $\omega_{\bt-1} = 0$.
The main difference is given by the fact that $\frac{1}{2}(a^{(L)} + b^{(L)})$  may a priori take values in an interval of width $1$.
As in the previous case we have $\alpha \le \delta_{\vb{\infty}}^{(L)} \le \alpha + \frac{1}{2}$ with $\alpha$ technically arbitrary.
We cannot thus choose a vanishing $k_{\delta^{(L)}_{\vb{\infty}}}$ but we have to consider $k_{\delta^{(L)}_{\infty}} = 0,\, 1$.

We find a third relation by considering the entry
\begin{equation}
  \Im\left(
    e^{+2\pi i \delta_{\vb{\infty}}^{(L)}}\,
    D^{(L)}\,
    \rM_{\vb{\infty}}^{(L)}\,
    \left( D^{(L)} \right)^{-1}
  \right)_{11}
  =
  \Im\left(
    \cL(n_{\vb{\infty}})
  \right)_{11}.
\end{equation}
Using
\begin{equation}
  \det \cC
  =
  \frac{\sin(\pi c^{(L)})}{\sin(\pi(a^{(L)}-b^{(L)}))},
\end{equation}
and the second equation in~\eqref{eq:cL} and~\eqref{eq:aL-bL} leads to:
\begin{equation}
  \cos(\pi( a^{(L)} + b^{(L)} - c^{(L)} ))
  =
  (-1)^{k_c+k_{\delta^{(L)}_{\vb{\infty}}} }\, \cos(2\pi \cA^{(L)}),
\end{equation}
where
\begin{equation}
  \cos(2\pi \cA^{(L)})
  =
  \cos(2\pi n_{\vb{0}})\,
  \cos(2\pi n_{\vb{\infty}})
  -
  \sin(2\pi n_{\vb{0}})\,
  \sin(2\pi n_{\vb{\infty}})\,
  \frac{n_{\vb{\infty}}^3}{n_{\vb{\infty}}}.
\label{eq:cos_n1}
\end{equation}
This expression is connected with rotation parameter in the third interaction point $\omega_{\bt+1} = 1$.
In fact $\cos(2\pi \cA^{(L)}) = \cos(2\pi {n}_{\vb{1}})$.
We then write
\begin{equation}
  a^{(L)} + b^{(L)} - c^{(L)}
  =
  2\, (-1)^{f^{(L)}}\, n_{\vb{1}}
  +
  k_c
  +
  k_{\delta^{(L)}_{\vb{\infty}}}
  +
  2\, k_{abc},
  \qquad
  k_{abc}\in \Z,
\end{equation}
with $f^{(L)} \in \left\lbrace 0, 1 \right\rbrace$.
The request
\begin{equation}
  A
  +
  B
  -
  n_{\vb{0}}
  -
  m_{\vb{0}}
  -
  (-1)^{f^{(L)}}\, n_{\vb{1}}
  -
  (-1)^{f^{(R)}}\, m_{\vb{1}}
  \in \Z
\end{equation}
finally fixes the $B$ parameter in the third equation of~\eqref{eq:parameters_equality_infty}.

So far we can summarise the results in
\begin{eqnarray}
  a
  =
  n_{\vb{0}} + (-1)^{f^{(L)}} n_{\vb{1}} + n_{\vb{\infty}} + m_a,
  & \qquad &
  m_a \in \Z,
  \\
  b
  =
  n_{\vb{0}} + (-1)^{f^{(L)}} n_{\vb{1}} - n_{\vb{\infty}} + m_b,
  & \qquad &
  m_b \in \Z,
  \\
  c
  =
  2\, n_{\vb{0}} + m_c,
  & \qquad &
  m_c \in \Z,
  \\
  \delta_{\vb{0}}^{(L)}
  =
  n_{\vb{0}},
  \\
  \delta_{\vb{\infty}}^{(L)}
  =
  - n_{\vb{0}} - (-1)^{f^{(L)}} n_{\vb{1}} + m_c + 2\, m_\delta,
  & \qquad &
  m_{\delta} \in \Z,
  \\
  A
  =
  n_{\vb{0}} + m_{\vb{0}} + m_A,
  & \qquad &
  m_A \in \Z,
  \\
  B
  =
  (-1)^{f^{(L)}}\, n_{\vb{1}} + (-1)^{f^{(R)}}\, m_{\vb{1}} + m_B,
  & \qquad &
  m_B \in \Z.
\end{eqnarray}

$K^{(L)}$ is finally determined from
\begin{equation}
  \left( D^{(L)}\, \rM_{\vb{\infty}}\, \left( D^{(L)} \right)^{-1} \right)_{21}
  =
  e^{-2\pi i \delta_{\vb{\infty}}^{(L)}}\,
  \left( \cL(n_{\vb{\infty}}) \right)_{21},
  \label{eq:fixing_K_21}
\end{equation}
and get:
\begin{equation}
  K^{(L)}
  =
  -\frac{(-1)^{m_a + m_b + m_c}}{2 \pi^2}\,
  \cG( a^{(L)},\, b^{(L)},\, c^{(L)} )\,
  \sin(2 \pi n_{\vb{0}})
  \sin(2 \pi n_{\vb{\infty}})
  \frac{n^1_{\vb{\infty}} + i\, n^2_{\vb{\infty}}}{n_{\vb{\infty}}},
  \label{eq:app_B_K21}
\end{equation}
where $\cG( a,\, b,\, c ) = \gfun{1-a}\, \gfun{1-b}\, \gfun{a+1-c}\, \gfun{b+1-c}$.


\subsection{Checking the Consistency of the Solution}

We check the consistency condition \eqref{eq:K_consistency_condition} using~\eqref{eq:product_in_SU2}.
The result is
\begin{equation}
  \begin{split}
    \left( K^{(L)} \right)^{-1}
    & =
    \frac{(-1)^{m_a + m_b + m_c}}{2 \pi^2}\,
    \cG(1 - a^{(L)},\, 1 - b^{(L)},\, 2 - c^{(L)})\,
    \\
    & \times
    \sin(2 \pi n_{\vb{0}})\,
    \sin(2 \pi n_{\vb{\infty}})\,
    \frac{n^1_{\vb{\infty}} -i n^2_{\vb{\infty}}}{n_{\vb{\infty}}},
  \end{split}
  \label{eq:app_B_K12}
\end{equation}
where the function $\cG( a,\, b,\, c )$ was defined at the end of the previous section.
Compatibility with~\eqref{eq:app_B_K21} requires
\begin{equation}
  \frac{(n^1_{\vb{\infty}})^2 + (n^2_{\vb{\infty}})^2}{n^2_{\vb{\infty}}}
  =
  -4 \frac{\sin(\pi a) \sin(\pi(c-a))\sin(\pi b) \sin(\pi(c-b))}
          {\sin^2(\pi c) \sin^2(\pi(a-b))}.
  \label{eq:n12+n22}
\end{equation}
We can then rewrite~\eqref{eq:cos_n1} as
\begin{equation}
  \frac{(n^3_{\vb{\infty}})^2}{n^2_{\vb{\infty}}}
  =
  \frac{(\cos(\pi (a-b)) \cos(\pi c)- \cos(\pi(a+b-c)))^2}
       {\sin^2(\pi c) \sin^2(\pi(a-b))}.
\end{equation}
It is then possible to verify that the sum of the left and right hand sides of~\eqref{eq:n12+n22} and the last equation are equal to $1$.
The same consistency check can also be performed by computing $K^{(L)}$ from
\begin{equation}
  \left( D^{(L)}\, \rM_{\vb{\infty}}\, \left( D^{(L)} \right)^{-1} \right)_{12}
  =
  e^{-2\pi i \delta_{\vb{\infty}}^{(L)}}\,
  \left( \cL(n_{\vb{\infty}}) \right)_{12},
\end{equation}
instead of \eqref{eq:fixing_K_21}.

% vim: ft=tex
